\documentclass[letterpaper,twocolumn,10pt]{article}
\usepackage{usenix-2020-09}
\usepackage{enumitem}
\setlist{nosep,leftmargin=*}

\begin{document}

\date{}
\title{\Large \bf KairOS: Adapting Wakeup Preemption to Load and Task Class\\with sched\_ext}
\author{{\rm Aayush Nair}}
\maketitle

\section{Motivation and Scheduling Problem}

Servers often run a latency-critical service next to batch jobs to keep
their CPUs busy~\cite{heracles,caladan}. Each time a request arrives, a
service thread wakes up. If every CPU is busy, the scheduler has to decide
at once: \emph{preempt} a running task so the request is served now, or let
the running task finish its timeslice while the request waits.

Neither fixed answer is always right. Preempting gives fast service when
the machine is quiet and costs a batch job little. But each preemption also
costs a context switch and pollutes caches and TLBs, which can cost more
than the switch itself~\cite{flexsc}, and when the CPUs are full of request
threads, preempting one to start another only moves the delay between
requests. On one core, Linux CFS preempted Apache's request injector
millions of times, while FreeBSD ULE, which does not preempt on wakeup, ran
Apache faster but can starve batch work~\cite{battle}.

Linux's current scheduler, EEVDF, preempts when a woken task has an earlier
virtual deadline and protects the running task until it has had its fair
share (\texttt{RUN\_TO\_PARITY})~\cite{eevdf}. Proposals to protect running
tasks further have been held back because they caused
starvation~\cite{r2pw}. Recent sched\_ext schedulers such as
\texttt{scx\_bpfland} and \texttt{scx\_lavd} shrink timeslices as load rises,
and \texttt{scx\_lavd} lets only its most latency-critical tasks
preempt~\cite{scx}, but none of them turns wakeup preemption on or off
according to measured load, the way Linux NAPI switches from interrupts to
polling under heavy network load.

\textbf{Problem.} For every wakeup, decide \emph{whether} to preempt and
\emph{whom}, from the current (and predicted) load and the kind of task that
is running, so that the service gets low tail latency both when the system
is quiet and during bursts, batch work keeps high throughput, and no task
starves.

\section{Proposed Policy}

KairOS is a scheduler written in C as a sched\_ext BPF
program~\cite{schedext}, with a small user-space loader. Every scheduling
decision is made in the kernel; user space only loads the policy, reads
statistics and writes parameters.

\textbf{Load signal.} The scheduler callbacks count runnable tasks, wakeups,
and the time tasks wait between becoming ready and running. A BPF timer
turns these into a fast and a slow moving average of runnable tasks per CPU
every millisecond (integer EWMAs). The wakeup rate is an early sign of a
burst; the queue wait confirms that it is hurting.

\textbf{Two modes.} In LATENCY mode, a woken task that finds no idle CPU may
preempt a busy CPU, and slices are short. In BATCH mode, only
latency-critical tasks within their budget may preempt, and batch tasks get
long, protected slices. The timer switches modes with asymmetric hysteresis:
KairOS returns to LATENCY mode as soon as the fast average drops below a low
threshold (\emph{fast exit}), but enters BATCH mode only when both averages
stay above a high threshold (\emph{cautious entry}). The thresholds are
calibrated by experiment.

\textbf{Class-aware preemption.} Whether a preemption helps depends on what
it interrupts. KairOS keeps cheap per-task features in BPF task storage:
how often the task goes to sleep rather than using up its slice, its average
run length, its CPU share, and who wakes whom. A small decision tree labels
each task latency-critical, normal, batch or background. A woken task may
only preempt a CPU running lower-class work, and then the task furthest
ahead in its fair share. ``Other tasks wait on it'' counts for more than
``it sleeps a lot''~\cite{critstacks}.

\textbf{Safety.} Latency-critical tasks get a CPU budget per window (e.g.,
70\% of 100\,ms); once it is used up, they lose their preemption rights
until the window resets. Every task also has a maximum wait (e.g., 50\,ms),
after which it runs next regardless of mode or class, well below the
sched\_ext watchdog timeout.

\textbf{Prediction (extension).} A user-space daemon reads the scheduler's
statistics every few seconds, predicts bursts, and writes a hint into a BPF
map. The kernel ignores stale hints, and a prediction may only move KairOS
into BATCH mode early; it can never delay the fast exit.

Forced always- and never-preempt modes of the same code serve as baselines
that differ from KairOS only in policy.

\section{Evaluation Plan}

\textbf{Setup.} Linux 7.0 with sched\_ext in a 4-vCPU VM; the load
generator runs on its own CPU.

\textbf{Workloads.} Latency-critical: a gRPC service (Go) whose worker
threads do a calibrated amount of CPU work per request, and Redis. Batch:
zstd compression (cache-heavy) and stress-ng. Load: an open-loop generator
with Poisson arrivals at steady rates, in synthetic on/off bursts of
different lengths, and from a scaled replay of the Azure Functions
trace~\cite{azure,shahrad}.

\textbf{Metrics.} Request latency (p50, p99, p99.9), batch throughput,
the longest wait of any task, preemptions, mode switches and reaction time,
and scheduler overhead.

\textbf{Baselines.} EEVDF (the Linux default), \texttt{scx\_bpfland},
\texttt{scx\_lavd} (default and performance mode), and KairOS forced to
always or never preempt. Ablations remove classification, the budget and
maximum wait, or prediction.

\begin{table}[h]
\centering
\small
\begin{tabular}{@{}ll@{}}
\hline
Experiment & Shows \\
\hline
E1 Steady-load sweep & Where fixed policies cross; thresholds \\
E2 Bursty load vs.\ baselines & Main result, per load phase \\
E3 Burst-length sweep & Burst length vs.\ reaction time \\
E4 Mode transitions & Latency spikes at switches \\
E5 Classification ablation & Value of class-aware preemption \\
E6 Starvation stress & Maximum-wait bound holds \\
E7 Overhead & Cost per callback \\
E8 Prediction & Prediction vs.\ simple persistence \\
\hline
\end{tabular}
\end{table}

\textbf{Method.} At least 10 interleaved runs per configuration; medians,
ranges and paired comparisons, against a measured noise floor.

\section{Expected Outcomes and Challenges}

\textbf{Expected outcomes.}
\begin{itemize}
\item Adaptive KairOS matches or beats the best fixed policy in each phase
  of a bursty workload: tail latency close to always-preempt when quiet and
  during bursts, batch throughput close to never-preempt.
\item Class-aware preemption removes most of the cost of preempting at high
  load, since request threads are not preempted to serve other requests.
\item No task waits longer than the maximum-wait bound.
\item An explanation of when adaptivity pays off (bursts longer than the
  reaction time) and when prediction helps (mainly periodic patterns).
\end{itemize}

\textbf{Challenges.}
\begin{itemize}
\item \emph{Noise.} Tail latency in a VM varies widely between identical
  runs.
\item \emph{Policy-dependent signals.} Preemption changes queue lengths,
  so the measured load moves when the mode changes.
\item \emph{Classification} must be cheap enough for BPF and stay right
  under overload, when busy request threads stop sleeping and look like
  batch work.
\item \emph{Runtimes that schedule too.} Go queues threads beyond
  \texttt{GOMAXPROCS} itself, out of the kernel's sight.
\item \emph{Starvation.} Protection rules can starve tasks, as the EEVDF
  proposals did; the budget and maximum wait must bound this within BPF's
  limits (verifier, integer arithmetic, watchdog).
\end{itemize}

{\footnotesize
\begin{thebibliography}{99}
\setlength{\itemsep}{0pt}

\bibitem{heracles}
D.~Lo et al.
\newblock Heracles: Improving resource efficiency at scale.
\newblock In {\em ISCA}, 2015.

\bibitem{caladan}
J.~Fried et al.
\newblock Caladan: Mitigating interference at microsecond timescales.
\newblock In {\em OSDI}, 2020.

\bibitem{flexsc}
L.~Soares and M.~Stumm.
\newblock FlexSC: Flexible system call scheduling with exception-less system
  calls.
\newblock In {\em OSDI}, 2010.

\bibitem{battle}
J.~Bouron et al.
\newblock The battle of the schedulers: FreeBSD ULE vs.\ Linux CFS.
\newblock In {\em USENIX ATC}, 2018.

\bibitem{eevdf}
Completing the EEVDF scheduler.
\newblock LWN.net, 2024. \url{https://lwn.net/Articles/969062}.

\bibitem{r2pw}
RUN\_TO\_PARITY\_WAKEUP patch discussion.
\newblock LKML, July 2026.
\newblock
  \url{https://patchew.org/linux/20260712113106.2564384-1-chenjinghuang2@huawei.com/}.

\bibitem{scx}
sched\_ext schedulers (\texttt{scx\_bpfland}, \texttt{scx\_lavd}).
\newblock \url{https://github.com/sched-ext/scx}.

\bibitem{schedext}
Extensible scheduler class (sched\_ext).
\newblock Linux kernel documentation.
\newblock \url{https://docs.kernel.org/scheduler/sched-ext.html}.

\bibitem{critstacks}
K.~Du~Bois et al.
\newblock Criticality stacks: Identifying critical threads in parallel
  programs using synchronization behavior.
\newblock In {\em ISCA}, 2013.

\bibitem{shahrad}
M.~Shahrad et al.
\newblock Serverless in the wild: Characterizing and optimizing the serverless
  workload at a large cloud provider.
\newblock In {\em USENIX ATC}, 2020.

\bibitem{azure}
Azure public dataset.
\newblock \url{https://github.com/Azure/AzurePublicDataset}.

\end{thebibliography}
}

\end{document}
